\section{核技巧：只通过内积看世界}
\label{sec:13-Machine-Learning/09-Kernel-Methods-and-Gaussian-Processes/02}

\hyperref[sec:13-Machine-Learning/09-Kernel-Methods-and-Gaussian-Processes/01]{上一节}的结论很直接：升到高维确实管用，但显式构造 \(\phi(x)\) 再在高维空间做内积，计算量追不上维度增长。这一节把上节末尾一闪而过的想法变成正式命题，然后处理两个更务实的问题：盘点哪些已学过的算法天然只通过内积接触数据；给``随手写一个相似度函数当核用''这件事提供一个判据。

核技巧的表述可以在三行内写完，但它的两个隐门槛才是出事的地方。第一个门槛：算法必须只以 \(\ip{x_i}{x_j}\) 的形式使用数据，不能出现 \(x_i\) 的孤立坐标。第二个门槛：放进算法的函数必须在任意有限样本上产生半正定 Gram 矩阵。两者缺一个，整个构造就不是``把算法搬进了某个特征空间''，而只是换了一个相似度度量——后者的几何后果不可控。

\subsection{核函数：不算特征，只算相似}

设有特征空间 \(\mathcal{F}\)（一个内积空间）与映射 \(\phi:\mathcal{X}\to\mathcal{F}\)。称

\[
k(x,z)=\ip{\phi(x)}{\phi(z)}_{\mathcal{F}}
\]

为由 \(\phi\) 诱导的核函数。核技巧的全部操作是：凡是算法中以 \(\ip{x}{z}\) 形式出现的项，换成 \(k(x,z)\)，整个算法就等价于在 \(\mathcal{F}\) 中原样运行，而 \(\phi\) 的坐标从头到尾不需要出现。

这不是近似，是恒等替换。代价则是丧失了在原空间中解释参数的能力——学到的权重向量变成了 \(\mathcal{F}\) 中的对象，我们只能通过 \(k\) 读取它与新样本的相似度。

\subsection{支持向量机：对偶里只有内积}

\hyperref[sec:09-Convex-Optimization/07-Applications/03]{``最大间隔与支持向量机''}中，硬间隔 SVM 的驻点条件给出 \(w=\sum_i\alpha_iy_ix_i\) 与 \(\sum_i\alpha_iy_i=0\)。把 \(w\) 的形式代回原始问题消去 \((w,b)\)，对偶问题只剩乘子：

\[
\max_{\alpha}\ \sum_{i=1}^{n}\alpha_i
-\frac12\sum_{i=1}^{n}\sum_{j=1}^{n}
\alpha_i\alpha_j\,y_iy_j\,\ip{x_i}{x_j},
\qquad
\alpha_i\ge0,\ \sum_{i=1}^{n}\alpha_iy_i=0.
\]

预测函数同样只用内积：

\[
f(x)=\sign\Bigl(\sum_{i=1}^{n}\alpha_iy_i\,\ip{x_i}{x}+b\Bigr).
\]

目标、约束、预测三处都只出现了 \(\ip{\cdot}{\cdot}\)。把每个 \(\ip{x_i}{x_j}\) 换成 \(k(x_i,x_j)\)，得到的是在特征空间里寻找最大间隔超平面；决策边界拉回原空间，则弯曲成 \(k\) 决定的形状。上一节的 XOR 例子：\(\phi(x)=(x_1,x_2,x_1x_2)\) 可分，对应核

\[
k(x,z)=x_1z_1+x_2z_2+x_1x_2\,z_1z_2.
\]

三维向量 \(\phi\) 不必出现，直接算核值即可。

\subsection{Ridge 回归：同一个解的 Gram 形态}

\hyperref[sec:13-Machine-Learning/02-Linear-Models/02]{``正则化从 Ridge 到 Lasso''}的一阶条件是

\[
(X^{\trans}X+\lambda I)\,w=X^{\trans}y.
\]

把解写成训练样本的线性组合 \(w=X^{\trans}\alpha\)（下一节的表示定理会给出这件事的一般理由），代入：

\[
X^{\trans}XX^{\trans}\alpha+\lambda X^{\trans}\alpha=X^{\trans}y.
\]

若

\[
(XX^{\trans}+\lambda I)\,\alpha=y
\]

成立，上式自然满足。记 Gram 矩阵 \(G=XX^{\trans}\)，\(G_{ij}=\ip{x_i}{x_j}\)，则

\[
\alpha=(G+\lambda I)^{-1}y,
\qquad
f(x)=\ip{w}{x}=\sum_{i=1}^{n}\alpha_i\,\ip{x_i}{x}.
\]

再一次只剩内积。这个形式还透露一个结构性事实：无论特征有多少维，\(w\) 永远躺在训练样本张成的空间里；对偶系数 \(\alpha\in\R^n\) 才是问题的真正自由度。把 \(G\) 换成核矩阵 \(K\)，就是本单元最后一节的核岭回归。

\subsection{PCA：协方差可以换成 Gram}

\hyperref[sec:13-Machine-Learning/04-Dimensionality-Reduction/01]{``PCA 从最大方差到最小重构误差''}中，主方向是 \(X_c^{\trans}X_c\) 的特征向量，\(X_c\) 是中心化数据矩阵。但 \(X_c^{\trans}X_c\in\R^{p\times p}\) 与 Gram 矩阵 \(X_cX_c^{\trans}\in\R^{n\times n}\) 共享全部非零特征值：若 \(u_j\) 是 \(X_cX_c^{\trans}\) 对应 \(\sigma_j^2\) 的单位特征向量，则 \(v_j=X_c^{\trans}u_j/\sigma_j\) 是 \(X_c^{\trans}X_c\) 对应同一特征值的单位特征向量。主方向与主成分得分可完全从 Gram 矩阵恢复。把 Gram 换成核矩阵，就是核 PCA 的起点。这里只指出这扇门，不走进去了。

\subsection{三个常见核}

\textbf{线性核。} \(k(x,z)=\ip{x}{z}\)，对应 \(\phi\) 为恒等映射。它是最基本的检验：核化后的算法必须退化回原来的线性算法。如果一个核方法在参数空间里的分析想不通，先退到线性核看看发生了什么。

\textbf{多项式核。} \(k(x,z)=(1+\ip{x}{z})^{p}\)，对应``全部不超过 \(p\) 次单项式''的特征空间。以 \(d=2,p=2\) 为例，\(x=(x_1,x_2)\)，\(z=(z_1,z_2)\)，直接展开：

\[
(1+\ip{x}{z})^{2}
=\bigl(1+x_1z_1+x_2z_2\bigr)^2
=1+2x_1z_1+2x_2z_2+x_1^2z_1^2+x_2^2z_2^2+2x_1x_2\,z_1z_2.
\]

取

\[
\phi(x)=\bigl(1,\ \sqrt2\,x_1,\ \sqrt2\,x_2,\ x_1^2,\ x_2^2,\ \sqrt2\,x_1x_2\bigr),
\]

逐项核对可见 \(\ip{\phi(x)}{\phi(z)}\) 与上式一致。核值等于六维内积并非魔术，只是组合数的记账：\((1+\ip{x}{z})^p\) 展开后的每一项天然是``\(x\) 的某个单项式乘 \(z\) 的同一单项式''，系数收入根号即可。计算量差了一个量级：左边一次内积加一次幂，\(O(d)\)；右边先构造约 \(d^p/p!\) 维向量再做内积。

\textbf{Gaussian/RBF 核。}

\[
k(x,z)=\exp\bigl(-\gamma\norm{x-z}^{2}\bigr),\qquad \gamma>0.
\]

把 \(\norm{x-z}^{2}=\norm{x}^{2}+\norm{z}^{2}-2\ip{x}{z}\) 代入，对 \(\exp(2\gamma\ip{x}{z})\) Taylor 展开，可见它对应无穷维的特征空间——全部单项式带着衰减系数出现，没有任何截断次数。直觉上，\(\gamma\) 控制相似度的衰减速度：\(\gamma\) 大时只把紧邻的点判为相似，学到的函数可以剧烈起伏；\(\gamma\) 小时远处的点也互相拉扯，函数被迫平滑。这个旋钮直接调节有效复杂度，最后一节还会回到它。

\subsection{什么时候可以自己设计核}

反过来问：凭直觉写一个相似度函数 \(s(x,z)\)，什么时候能当真当核用？答案是正定性判据（机器学习语境中常被宽泛地称为 Mercer 条件；经典 Mercer 定理还包含定义域、连续性与测度等附加要求）：

\[
k\ \text{是某个特征空间的内积}
\iff
\text{对任意有限点集 } \{x_1,\dots,x_n\},
\ K=\bigl[k(x_i,x_j)\bigr]_{i,j}\ \text{半正定}.
\]

必要性只需一行：若 \(k(x,z)=\ip{\phi(x)}{\phi(z)}\)，则对任意 \(v\in\R^n\)，

\[
v^{\trans}Kv
=\sum_{i=1}^{n}\sum_{j=1}^{n}v_iv_j\,\ip{\phi(x_i)}{\phi(x_j)}
=\norm{\sum_{i=1}^{n}v_i\,\phi(x_i)}^{2}\ge0.
\]

充分性——从半正定二元函数真的构造出一个特征空间——是下一节构造 RKHS 的内容。

这个判据让``设计核''从玄学变成可检验的工程。核的封闭性提供了积木：两个核的和仍是核，两个核的积仍是核，核乘正常数仍是核——对应 Gram 矩阵的加、逐元素乘、标量乘分别保持半正定。于是可以按数据的不同侧面分别构造相似度再组合。反过来，凭直觉写出却不满足半正定的``相似度''，会让 SVM 对偶目标失去凹性、让``范数''不再是范数——\hyperref[sec:09-Convex-Optimization/07-Applications/03]{``最大间隔与支持向量机''}末尾对非 PSD 相似度的警告正来自这层根源。

\subsection{核技巧没回答的问题}

核技巧至此仍只是算法层面的替换规则。它没有回答：\(k\) 对应的特征空间到底长什么样，里面的``函数''是什么数学对象，\(\norm{w}^{2}\) 换成核语言后被惩罚的究竟是什么。不回答这些问题，``加正则项控制复杂度''就是一句喊起来对但不知对在哪里的口号。

\hyperref[sec:13-Machine-Learning/09-Kernel-Methods-and-Gaussian-Processes/03]{``再生核希尔伯特空间：把函数当成点''}把这件事严格回答清楚。学完那篇再回头看这里盘点的三个算法核化版本，会发现它们不是三个孤立的技巧，而是同一个 RKHS 正则化框架在不同损失函数下的三个实例。
